\chapter[\texorpdfstring{Достаточные условия дифференцируемости функции не\-скольких переменных.}{Достаточные условия дифференцируемости функции нескольких переменных}]{Достаточные условия дифференцируемости функции нескольких переменных.}
\section[\texorpdfstring{Достаточные условия дифференцируемости функции \\несколь\-ких переменных.}{Достаточные условия дифференцируемости функции нескольких переменных}]{Достаточные условия дифференцируемости функции нескольких переменных}

%----------------------------------------------------------------------------------------
% Определение 10.1 --- Частные производные
%----------------------------------------------------------------------------------------

\begin{defn} \label{df10:1}
Частной производной по $x$ функции двух переменных $f(x, y)$ в точке $(x_0, y_0)$ называется производная функции $f(x, y_0)$ одной переменной $x$ в точке $x_0$ (обозначается $\frac{\partial f}{\partial x}(x_0, y_0)$ или $f'_x(x_0, y_0)$; при этом предполагается, что функция $f(x, y_0)$ определена в некоторой окрестности точки $x_0$ и её производная в точке $x_0$ конечна). Частной производной по $y$ функции двух переменных $f(x, y)$ в точке $(x_0, y_0)$ называется производная функции $f(x_0, y)$ одной переменной $y$ в точке $y_0$ (обозначается $\frac{\partial f}{\partial y}(x_0, y_0)$ или $f'_y(x_0, y_0)$; при этом предполагается, что функция $f(x_0, y)$ определена в некоторой окрестности точки $y_0$ и её производная в точке $y_0$ конечна).

Эти определения могут быть записаны так:
$$
\frac{\partial f}{\partial x}(x_0, y_0) = \frac{d}{dx} f(x, y_0) \bigg|_{x=x_0};
$$
$$
\frac{\partial f}{\partial y}(x_0, y_0) = \frac{d}{dy} f(x_0, y) \bigg|_{y=y_0}.
$$
\end{defn}

%----------------------------------------------------------------------------------------
% Определение 10.2 --- Дифференцируемость функции многих переменных
%----------------------------------------------------------------------------------------

\begin{defn} \label{df10:2}
Функция $n$ переменных $f$ называется дифференцируемой в точке $(x_1^0, \ldots, x_n^0)$, если её приращение в этой точке может быть представлено в виде
\begin{multline} \label{eq10:1}
\Delta f(x_1^0, \ldots, x_n^0) \equiv f(x_1^0 + \Delta x_1, \ldots, x_n^0 + \Delta x_n) - f(x_1^0, \ldots, x_n^0) = \\
= A_1\Delta x_1 + \ldots + A_n\Delta x_n + \alpha(\Delta x_1, \ldots, \Delta x_n) \times \\
\times \sqrt{(\Delta x_1)^2 + \ldots + (\Delta x_n)^2},
\end{multline}
где $A_i \in \bbR$, $i = 1, 2, \ldots, n$, а предел функции $n$ переменных $\alpha(\Delta x_1, \ldots, \Delta x_n)$ равен $0$ при $(\Delta x_1, \ldots, \Delta x_n) \to (0, 0, \ldots, 0)$. При этом линейная часть приращения $A_1\Delta x_1 + \ldots + A_n\Delta x_n$ называется дифференциалом функции $f$ в точке $(x_1^0, \ldots, x_n^0)$ и обозначается $df(x_1^0, \ldots, x_n^0)$.

Иначе равенство \eqref{eq10:1} можно записать так:
$$
\Delta f(x_1^0, \ldots, x_n^0) = A_1\Delta x_1 + \ldots + A_n\Delta x_n + o(\rho)
$$
при $(\Delta x_1, \ldots, \Delta x_n) \to (0, \ldots, 0)$, где
$$
\rho = \sqrt{(\Delta x_1)^2 + \ldots + (\Delta x_n)^2}.
$$
Как и в одномерном случае, $o(\varphi(x_1, \ldots, x_n))$ понимается как $\alpha(x_1, \ldots, x_n) \cdot \varphi(x_1, \ldots, x_n)$, где предел функции $\alpha(x_1, \ldots, x_n)$ равен $0$ при соответствующем стремлении $x_1, \ldots, x_n$.

Для функции двух переменных равенство \eqref{eq10:1} запишется так:
$$
f(x_0 + \Delta x, y_0 + \Delta y) - f(x_0, y_0) =
$$
$$
= A\Delta x + B\Delta y + o\left(\sqrt{(\Delta x)^2 + (\Delta y)^2}\right),
$$
т.е. $\Delta f(x_0, y_0) = df(x_0, y_0) + o(\rho)$, при $(\Delta x, \Delta y) \to (0, 0)$.
\end{defn}

%----------------------------------------------------------------------------------------
% Теорема 10.1 --- Необходимые условия дифференцируемости
%----------------------------------------------------------------------------------------

\begin{thm}[необходимые условия дифференцируемости] \label{th10:1}
Если функция $n$ переменных $f$ дифференцируема в точке $(x_1^0, \ldots, x_n^0)$, то:
\begin{enumerate}
    \item[1)] $f$ непрерывна в точке $(x_1^0, \ldots, x_n^0)$;
    \item[2)] при $i = 1, 2, \ldots, n$ существуют конечные $\frac{\partial f}{\partial x_i}(x_1^0, \ldots, x_n^0)$, соответственно равные коэффициентам $A_i$ в линейной части приращения.
\end{enumerate}
\end{thm}

Напомним, что для функций одной переменной дифференцируемость равносильна существованию конечной производной в рассматриваемой точке. В случае двух переменных существование частных производных не является достаточным даже для непрерывности функции. Например, функция $f(x,y)=\sign xy $ разрывна в точке $(0, 0)$, однако в этой точке она имеет обе частные производные $f'_x$, $f'_y$, равные нулю. Поэтому рассмотрим достаточное условие дифференцируемости для функций двух переменных.
\begin{thm}[Достаточное условие дифференцируемости] \label{ch8th1}
Если функция $f(x,y)$ имеет в некоторой окрестности $O(x^0, y^0)$ частные производные $f'_x(x,y)$, $f'_y(x,y)$, $(x,y) \in O(x^0, y^0)$, которые непрерывны в точке $(x^0, y^0)$, то функция $f$ дифференцируема в точке $(x^0, y^0)$.
\end{thm}
\begin{proof}
Через $O_\delta(x^0, y^0)$ обозначим $\delta$-окрестность точки $(x^0, y^0)$, в которой существуют частные производные $f'_x(x,y)$, $f'_y(x,y)$. Пусть $(x,y) \in O_\delta(x^0, y^0)$, что по определению означает:
$$
|x - x^0|^2 + |y - y^0|^2 < \delta^2.
$$ 
Т.к. в этом неравенстве первое слагаемое неотрицательно, то отсюда и из геометрических соображений рисунка следует $|y - y^0|^2 < \delta^2$. Другими словами, точка $(x^0, y)$ принадлежит $O_\delta(x^0, y^0)$ и имеет смысл записать следующее равенство:
\begin{equation*}
f(x,y) - f(x^0, y^0) = (f(x,y) - f(x^0, y)) + (f(x^0, y) - f(x^0, y^0)),
\end{equation*}

Покажем, что к разностям в скобках можно применить теорему Лагранжа. Действительно, согласно условию теоремы функция $f(x, y)$ имеет частную производную $f'_x$ в любой точке $(\xi, y)$  при $\xi \in [x^0, x]$, т.к. тогда $(\xi, y) \in O_\delta(x^0, y^0)$ Проще говоря, функция $\phi(\xi)=f(\xi, y)$ дифференцируема на отрезке $[x^0, x]$ и к ней применима теорема Лагранжа: существует $\xi_1 \in (x^0, x)$, такое что выполняется $\phi(x) - \phi(x^0) = \phi'(\xi_1)(x - x^0)$. Или же, другими словами, $\exists \xi_1 \in (x^0, x)$:
$$
f(x,y) - f(x^0, y) = f'_x(\xi_1, y)(x - x^0)
$$

Аналогично, $\exists \eta_1 \in (y, y^0)$:
$$
f(x^0, y) - f(x^0, y^0) = f'_y(x^0, \eta_1)(y - y^0)
$$ 

Таким образом, показано, что для всех точек $(x,y)$ из окрестности $O_\delta(x^0, y^0)$ найдутся такие $\xi_1 \in (x^0, x)$ и $\eta_1 \in (y^0, y)$, при которых:
$$
f(x,y) - f(x^0, y^0) = f'_x(\xi_1, y) \Delta x + f'_y(x, \eta_1) \Delta y
$$

Т.к.~частные производные $f'_x$, $f'_y$ непрерывны в точке $(x^0, y^0)$, то 
$$
\lim\limits_{(x,y) \to (x^0, y^0)} f'_x(\xi_1, y) = f'_x(x^0, y^0);\qquad 
\lim\limits_{(x,y) \to (x^0, y^0)} f'_y(x, \eta_1) = f'_y(x^0, y^0).
$$
Следовательно, функции $\alpha =  f'_x(\xi_1, y) - f'_x(x^0, y^0)$ и $\beta =  f'_y(x, \eta_1) - f'_y(x^0, y^0)$ являются бесконечно малыми при $(x,y) \to (x^0, y^0)$. 

Таким образом,
$$
f(x,y) - f(x^0, y^0) = f'_x(x^0, y^0) \Delta x + f'_y(x^0, y^0) \Delta y + \alpha \Delta x + \beta \Delta y,
$$
где $\alpha \to 0$ и $\beta \to 0$ при $(x,y)\to(x^0,y^0)$.

Осталось лишь показать, что последние два слагаемые в сумме тоже представляют из себя бесконечно малую величину. Представим ее в следующем виде:
$$
\alpha \Delta x + \beta \Delta y = \frac{\alpha \Delta x + \beta \Delta y}{\sqrt{\Delta x^2 + \Delta y^2}} \cdot \sqrt{\Delta x^2 + \Delta y^2} = \gamma \cdot \sqrt{\Delta x^2 + \Delta y^2},
$$
где
$$
\gamma = \alpha \frac{\Delta x}{\sqrt{\Delta x^2 + \Delta y^2}} + \beta \frac{\Delta y}{\sqrt{\Delta x^2 + \Delta y^2}}.
$$
Поскольку $\alpha$ и $\beta$ "--- бесконечно малые, а соответствующие множители являются ограниченными функциями, то следовательно функция $\gamma$ сама является бесконечно малой и $\gamma \cdot \sqrt{\Delta x^2 + \Delta y^2} = o(\sqrt{\Delta x^2 + \Delta y^2})$. Окончательно
$$
f(x,y) - f(x^0, y^0) = f'_x(x^0, y^0) \Delta x + f'_y(x^0, y^0) \Delta y + o(\sqrt{\Delta x^2 + \Delta y^2}),
$$
т.е.~$f$ дифференцируема в точке $(x^0, y^0)$.

Теорема доказана.
\end{proof}

На самом деле, эта теорема переносится на случай многомерного пространства.
\begin{thmn}
Если функция $f(x)$ имеет частные производные $\partial f /\partial x_i$, $i=\overline{1,n}$, в некоторой окрестности точки $x_0$ и эти производные непрерывны в точке $x_0$, то функция $f(x)$ дифференцируема в точке $x_0$.
\end{thmn}
\begin{cons}
Если функция $f$ определена на области $G \subset \bbR^n$ и имеет на этой области частные производные 1-го порядка, которые непрерывны на области $G$, то функция $f$ дифференцируема в каждой точке области~$G$.
\end{cons}
